\section{Discussion}
\label{sec:discussion}

\subsection{What the evidence supports}

Correctness-gated policy optimization changes the one-draw distribution of RTL
quality.  The central evidence is the equal-sample endpoint: it preserves the
original probability mass, assigns failures zero, uses real post-route
measurements, and improves on \claim{OverallImproved} of
\claim{DesignCount} frozen designs.  The best-of-$N$ analysis further shows that
the change is practically meaningful for sampling efficiency.

The mechanism is best described as probability reallocation.  Supervised
training teaches multiple implementation styles; policy optimization raises the
mass on styles that the gate accepts and the proxy ranks highly.  This account
predicts smaller gains when the supervised support lacks a fast reliable style,
matching the weaker families in Table~\ref{tab:family-verified}.  We therefore do
not claim autonomous invention of unseen microarchitectures.

\subsection{Threats to validity}

\paragraph{Template-bounded generalization}
The frozen split establishes transfer to unseen parameters and extrapolation
beyond training ranges within the generator's design system.  It does not test
free-form industrial specifications, control-dominated systems, memories, or
multiple clock domains.

\paragraph{One device and one main training run}
Physical measurements target one FPGA family, and the primary policy result
comes from one training trajectory.  The paired design breadth makes the
evaluation stronger than a few selected examples, but independent training
seeds and another device would materially strengthen a journal version.

\paragraph{Timing-derived frequency}
The primary endpoint is derived from worst negative slack after implementation,
not from repeated timing closure at every candidate clock.  Both policies use
the same tool flow.  The symmetric board sweep is the planned external check,
and no historical asymmetric board value is used as evidence in this draft.

\paragraph{Post-hoc uncertainty}
The bootstrap intervals quantify sensitivity to the held-out design sample.
They were added after the artifacts existed and should not be interpreted as a
preregistered confirmatory test.

\subsection{Why the trajectory is a subplot}

The trajectory reveals an important operational failure, but it is downstream
of the main capability result.  Early training yields a large real improvement;
later training inflates proxy reward while measured conditional frequency is
flat and correctness falls.  This motivates periodic real-EDA audits and early
stopping criteria.  It does not justify replacing the paper with a general
reward-model-repair claim, especially because later architecture experiments
showed that the observed collapse was not universal.
